% Modernes Bildquellenverzeichnis
% Nutzt biblatex für Quellenangaben von Abbildungen

% === Sauberer Seitenumbruch ===
\clearpage
\cleardoublepage

% === Spezielle Formatierung für Quellenverzeichnis ===
\begingroup
\let\cleardoublepage\relax
\let\clearpage\relax

% === Institution Logo ===
\begin{figure}[t]
  \vspace{1.5cm}
  \begin{flushright}
    \IfFileExists{default_style/AtheneGespiegelt.png}{%
      \includegraphics[width=0.1\textwidth]{default_style/AtheneGespiegelt.png}%
    }{%
      \textcolor{secondary}{\rule{0.1\textwidth}{0.1\textwidth}}%
    }
  \end{flushright}
  \vspace{-4.5cm}
\end{figure}

% === Bildquellenverzeichnis drucken ===
\printbibliography[
  category=abbildung,
  title={Bildquellenverzeichnis}
]

\endgroup

% === Seitenfüllung und Abschluss ===
\vspace*{\fill}
\clearpage
\thispagestyle{empty}